% ============================================================================
% PROPOSED RESULTS BLOCK for mechanisms.tex — the scenario-impact story.
%
% Ownership: mechanisms.tex is session A's artifact. This file is the authoring
% session's block proposal (charter v2, lane 3): integrate by % ============================================================================
% PROPOSED RESULTS BLOCK for mechanisms.tex — the scenario-impact story.
%
% Ownership: mechanisms.tex is session A's artifact. This file is the authoring
% session's block proposal (charter v2, lane 3): integrate by % ============================================================================
% PROPOSED RESULTS BLOCK for mechanisms.tex — the scenario-impact story.
%
% Ownership: mechanisms.tex is session A's artifact. This file is the authoring
% session's block proposal (charter v2, lane 3): integrate by % ============================================================================
% PROPOSED RESULTS BLOCK for mechanisms.tex — the scenario-impact story.
%
% Ownership: mechanisms.tex is session A's artifact. This file is the authoring
% session's block proposal (charter v2, lane 3): integrate by \input{results-block}
% after the channel frames (before any conclusion frame), or copy frames in and
% delete this file. It uses only macros mechanisms.tex already defines
% (\figframe, \srcnote) and figures under figures/impact/ (committed on the
% writing branch; regenerate from the run record with the impact-chart script
% against any (case, base, reform, coupled) quadruple — one command on the next calibration).
%
% Register: conceptual flows only, no equations (Marcelo, 2026-08-13). Every
% chart carries its evidence base on-face (PEP vs Base, Philippines).
% ============================================================================

% -- transition frame: from mechanism to application -------------------------
\begin{frame}[plain]
  \vfill
  {\usebeamerfont{title}\Large\color{ink}What the coupled system says about
   one real programme}\par\vspace{6pt}
  {\small\color{sub}The Philippine decarbonisation composition (coal moratorium
   from 2027), run end to end:\\ energy system, emissions, health, and the
   economy --- one scenario, every consequence on its own account.}\par
  \vfill
  \srcnote{ogclews\_runs\_v16/coupled/ (manifest + macro\_table + figures);
    CLEWs Philippines\_v16\_CALIBRATED, Base\_v16/PEP\_v16. All charts name the
    base on-face; regenerated by experiments/scenario\_impact\_charts.py.}
\end{frame}

\figframe{The scenario rebuilds the power system}%
  {impact/impact_generation_mix}%
  {Generation by technology under the reform. Demand roughly quadruples by 2053;
   coal exits after the 2027 moratorium; offshore wind carries the expansion to
   its binding build limit. PEP vs Base, Philippines.}

\figframe{What gets built, and when}%
  {impact/impact_new_capacity}%
  {New capacity by era. The near decade is solar-led; offshore wind dominates
   the 2030s; the long tail adds onshore wind and gas peaking. PEP vs Base, Philippines.}

\figframe{Where the carbon comes from --- and what the scenario avoids}%
  {impact/impact_co2e}%
  {National CO2e by source, including land-conversion carbon (newly accounted;
   the sawtooth is conversion pulses). The cumulative avoided total vs the
   no-policy path is printed on the chart. PEP vs Base, Philippines.}

\figframe{Cleaner air, counted in lives and working time}%
  {impact/impact_air_health}%
  {The particulate change is converted through Global Burden of Disease
   dose--response and age profiles: avoided deaths (mostly at older ages) and
   returned working time (morbidity). Aggregate GDP effect below reporting
   precision. PEP vs Base, Philippines.}

\figframe{The economy pays modestly, and mostly through wages}%
  {impact/impact_macro}%
  {GDP, consumption, and wages against the no-policy baseline over the
   transition window; steady-state values labelled. The wage carries the
   adjustment; employment barely moves. PEP vs Base, Philippines.}

% -- the reading rules for the headline (no equations; the honesty slide) ----
\begin{frame}{Read the macro headline with its transmission named}
  \begin{itemize}\setlength\itemsep{8pt}
    \item Nearly all of the output effect arrives through one assumption:
          higher system costs pass into industry production costs
          (a reduced-form cost-push mapping, not a factor-demand system).
    \item Routing the \emph{same} cost signal through sectoral productivity
          instead \textbf{reverses the sign} of the aggregate response
          --- so any quoted headline names its transmission.
    \item Inside the applied bundle, carbon is \emph{emitted only} (no carbon
          price acts in these numbers), and investment and health contribute
          $\approx 0$ to the aggregate --- the coupled response is the energy
          composite by construction --- an accounting identity of the
          configuration, not a zero-interaction finding.
  \end{itemize}
  \srcnote{ogclews\_manifest.json channels[]/provenance[] (emit-only carbon,
    near-zero investment and health); the matched composition battery and the
    transmission sign-flip re-run are pending on the current calibration ---
    quote neither until they land.}
\end{frame}

\figframe{Who gains and who loses is not a budget-share story}%
  {impact/impact_incidence}%
  {Consumption-equivalent incidence by lifetime-income group on the coupled
   run: middle groups gain, the 90th--99th percentile group loses most, and
   baseline energy budget shares are nearly uniform --- the dispersion runs
   through earnings, assets, and transfers. PEP vs Base, Philippines.}

% after the channel frames (before any conclusion frame), or copy frames in and
% delete this file. It uses only macros mechanisms.tex already defines
% (\figframe, \srcnote) and figures under figures/impact/ (committed on the
% writing branch; regenerate from the run record with the impact-chart script
% against any (case, base, reform, coupled) quadruple — one command on the next calibration).
%
% Register: conceptual flows only, no equations (Marcelo, 2026-08-13). Every
% chart carries its evidence base on-face (PEP vs Base, Philippines).
% ============================================================================

% -- transition frame: from mechanism to application -------------------------
\begin{frame}[plain]
  \vfill
  {\usebeamerfont{title}\Large\color{ink}What the coupled system says about
   one real programme}\par\vspace{6pt}
  {\small\color{sub}The Philippine decarbonisation composition (coal moratorium
   from 2027), run end to end:\\ energy system, emissions, health, and the
   economy --- one scenario, every consequence on its own account.}\par
  \vfill
  \srcnote{ogclews\_runs\_v16/coupled/ (manifest + macro\_table + figures);
    CLEWs Philippines\_v16\_CALIBRATED, Base\_v16/PEP\_v16. All charts name the
    base on-face; regenerated by experiments/scenario\_impact\_charts.py.}
\end{frame}

\figframe{The scenario rebuilds the power system}%
  {impact/impact_generation_mix}%
  {Generation by technology under the reform. Demand roughly quadruples by 2053;
   coal exits after the 2027 moratorium; offshore wind carries the expansion to
   its binding build limit. PEP vs Base, Philippines.}

\figframe{What gets built, and when}%
  {impact/impact_new_capacity}%
  {New capacity by era. The near decade is solar-led; offshore wind dominates
   the 2030s; the long tail adds onshore wind and gas peaking. PEP vs Base, Philippines.}

\figframe{Where the carbon comes from --- and what the scenario avoids}%
  {impact/impact_co2e}%
  {National CO2e by source, including land-conversion carbon (newly accounted;
   the sawtooth is conversion pulses). The cumulative avoided total vs the
   no-policy path is printed on the chart. PEP vs Base, Philippines.}

\figframe{Cleaner air, counted in lives and working time}%
  {impact/impact_air_health}%
  {The particulate change is converted through Global Burden of Disease
   dose--response and age profiles: avoided deaths (mostly at older ages) and
   returned working time (morbidity). Aggregate GDP effect below reporting
   precision. PEP vs Base, Philippines.}

\figframe{The economy pays modestly, and mostly through wages}%
  {impact/impact_macro}%
  {GDP, consumption, and wages against the no-policy baseline over the
   transition window; steady-state values labelled. The wage carries the
   adjustment; employment barely moves. PEP vs Base, Philippines.}

% -- the reading rules for the headline (no equations; the honesty slide) ----
\begin{frame}{Read the macro headline with its transmission named}
  \begin{itemize}\setlength\itemsep{8pt}
    \item Nearly all of the output effect arrives through one assumption:
          higher system costs pass into industry production costs
          (a reduced-form cost-push mapping, not a factor-demand system).
    \item Routing the \emph{same} cost signal through sectoral productivity
          instead \textbf{reverses the sign} of the aggregate response
          --- so any quoted headline names its transmission.
    \item Inside the applied bundle, carbon is \emph{emitted only} (no carbon
          price acts in these numbers), and investment and health contribute
          $\approx 0$ to the aggregate --- the coupled response is the energy
          composite by construction --- an accounting identity of the
          configuration, not a zero-interaction finding.
  \end{itemize}
  \srcnote{ogclews\_manifest.json channels[]/provenance[] (emit-only carbon,
    near-zero investment and health); the matched composition battery and the
    transmission sign-flip re-run are pending on the current calibration ---
    quote neither until they land.}
\end{frame}

\figframe{Who gains and who loses is not a budget-share story}%
  {impact/impact_incidence}%
  {Consumption-equivalent incidence by lifetime-income group on the coupled
   run: middle groups gain, the 90th--99th percentile group loses most, and
   baseline energy budget shares are nearly uniform --- the dispersion runs
   through earnings, assets, and transfers. PEP vs Base, Philippines.}

% after the channel frames (before any conclusion frame), or copy frames in and
% delete this file. It uses only macros mechanisms.tex already defines
% (\figframe, \srcnote) and figures under figures/impact/ (committed on the
% writing branch; regenerate from the run record with the impact-chart script
% against any (case, base, reform, coupled) quadruple — one command on the next calibration).
%
% Register: conceptual flows only, no equations (Marcelo, 2026-08-13). Every
% chart carries its evidence base on-face (PEP vs Base, Philippines).
% ============================================================================

% -- transition frame: from mechanism to application -------------------------
\begin{frame}[plain]
  \vfill
  {\usebeamerfont{title}\Large\color{ink}What the coupled system says about
   one real programme}\par\vspace{6pt}
  {\small\color{sub}The Philippine decarbonisation composition (coal moratorium
   from 2027), run end to end:\\ energy system, emissions, health, and the
   economy --- one scenario, every consequence on its own account.}\par
  \vfill
  \srcnote{ogclews\_runs\_v16/coupled/ (manifest + macro\_table + figures);
    CLEWs Philippines\_v16\_CALIBRATED, Base\_v16/PEP\_v16. All charts name the
    base on-face; regenerated by experiments/scenario\_impact\_charts.py.}
\end{frame}

\figframe{The scenario rebuilds the power system}%
  {impact/impact_generation_mix}%
  {Generation by technology under the reform. Demand roughly quadruples by 2053;
   coal exits after the 2027 moratorium; offshore wind carries the expansion to
   its binding build limit. PEP vs Base, Philippines.}

\figframe{What gets built, and when}%
  {impact/impact_new_capacity}%
  {New capacity by era. The near decade is solar-led; offshore wind dominates
   the 2030s; the long tail adds onshore wind and gas peaking. PEP vs Base, Philippines.}

\figframe{Where the carbon comes from --- and what the scenario avoids}%
  {impact/impact_co2e}%
  {National CO2e by source, including land-conversion carbon (newly accounted;
   the sawtooth is conversion pulses). The cumulative avoided total vs the
   no-policy path is printed on the chart. PEP vs Base, Philippines.}

\figframe{Cleaner air, counted in lives and working time}%
  {impact/impact_air_health}%
  {The particulate change is converted through Global Burden of Disease
   dose--response and age profiles: avoided deaths (mostly at older ages) and
   returned working time (morbidity). Aggregate GDP effect below reporting
   precision. PEP vs Base, Philippines.}

\figframe{The economy pays modestly, and mostly through wages}%
  {impact/impact_macro}%
  {GDP, consumption, and wages against the no-policy baseline over the
   transition window; steady-state values labelled. The wage carries the
   adjustment; employment barely moves. PEP vs Base, Philippines.}

% -- the reading rules for the headline (no equations; the honesty slide) ----
\begin{frame}{Read the macro headline with its transmission named}
  \begin{itemize}\setlength\itemsep{8pt}
    \item Nearly all of the output effect arrives through one assumption:
          higher system costs pass into industry production costs
          (a reduced-form cost-push mapping, not a factor-demand system).
    \item Routing the \emph{same} cost signal through sectoral productivity
          instead \textbf{reverses the sign} of the aggregate response
          --- so any quoted headline names its transmission.
    \item Inside the applied bundle, carbon is \emph{emitted only} (no carbon
          price acts in these numbers), and investment and health contribute
          $\approx 0$ to the aggregate --- the coupled response is the energy
          composite by construction --- an accounting identity of the
          configuration, not a zero-interaction finding.
  \end{itemize}
  \srcnote{ogclews\_manifest.json channels[]/provenance[] (emit-only carbon,
    near-zero investment and health); the matched composition battery and the
    transmission sign-flip re-run are pending on the current calibration ---
    quote neither until they land.}
\end{frame}

\figframe{Who gains and who loses is not a budget-share story}%
  {impact/impact_incidence}%
  {Consumption-equivalent incidence by lifetime-income group on the coupled
   run: middle groups gain, the 90th--99th percentile group loses most, and
   baseline energy budget shares are nearly uniform --- the dispersion runs
   through earnings, assets, and transfers. PEP vs Base, Philippines.}

% after the channel frames (before any conclusion frame), or copy frames in and
% delete this file. It uses only macros mechanisms.tex already defines
% (\figframe, \srcnote) and figures under figures/impact/ (committed on the
% writing branch; regenerate from the run record with the impact-chart script
% against any (case, base, reform, coupled) quadruple — one command on the next calibration).
%
% Register: conceptual flows only, no equations (Marcelo, 2026-08-13). Every
% chart carries its evidence base on-face (PEP vs Base, Philippines).
% ============================================================================

% -- transition frame: from mechanism to application -------------------------
\begin{frame}[plain]
  \vfill
  {\usebeamerfont{title}\Large\color{ink}What the coupled system says about
   one real programme}\par\vspace{6pt}
  {\small\color{sub}The Philippine decarbonisation composition (coal moratorium
   from 2027), run end to end:\\ energy system, emissions, health, and the
   economy --- one scenario, every consequence on its own account.}\par
  \vfill
  \srcnote{ogclews\_runs\_v16/coupled/ (manifest + macro\_table + figures);
    CLEWs Philippines\_v16\_CALIBRATED, Base\_v16/PEP\_v16. All charts name the
    base on-face; regenerated by experiments/scenario\_impact\_charts.py.}
\end{frame}

\figframe{The scenario rebuilds the power system}%
  {impact/impact_generation_mix}%
  {Generation by technology under the reform. Demand roughly quadruples by 2053;
   coal exits after the 2027 moratorium; offshore wind carries the expansion to
   its binding build limit. PEP vs Base, Philippines.}

\figframe{What gets built, and when}%
  {impact/impact_new_capacity}%
  {New capacity by era. The near decade is solar-led; offshore wind dominates
   the 2030s; the long tail adds onshore wind and gas peaking. PEP vs Base, Philippines.}

\figframe{Where the carbon comes from --- and what the scenario avoids}%
  {impact/impact_co2e}%
  {National CO2e by source, including land-conversion carbon (newly accounted;
   the sawtooth is conversion pulses). The cumulative avoided total vs the
   no-policy path is printed on the chart. PEP vs Base, Philippines.}

\figframe{Cleaner air, counted in lives and working time}%
  {impact/impact_air_health}%
  {The particulate change is converted through Global Burden of Disease
   dose--response and age profiles: avoided deaths (mostly at older ages) and
   returned working time (morbidity). Aggregate GDP effect below reporting
   precision. PEP vs Base, Philippines.}

\figframe{The economy pays modestly, and mostly through wages}%
  {impact/impact_macro}%
  {GDP, consumption, and wages against the no-policy baseline over the
   transition window; steady-state values labelled. The wage carries the
   adjustment; employment barely moves. PEP vs Base, Philippines.}

% -- the reading rules for the headline (no equations; the honesty slide) ----
\begin{frame}{Read the macro headline with its transmission named}
  \begin{itemize}\setlength\itemsep{8pt}
    \item Nearly all of the output effect arrives through one assumption:
          higher system costs pass into industry production costs
          (a reduced-form cost-push mapping, not a factor-demand system).
    \item Routing the \emph{same} cost signal through sectoral productivity
          instead \textbf{reverses the sign} of the aggregate response
          --- so any quoted headline names its transmission.
    \item Inside the applied bundle, carbon is \emph{emitted only} (no carbon
          price acts in these numbers), and investment and health contribute
          $\approx 0$ to the aggregate --- the coupled response is the energy
          composite by construction --- an accounting identity of the
          configuration, not a zero-interaction finding.
  \end{itemize}
  \srcnote{ogclews\_manifest.json channels[]/provenance[] (emit-only carbon,
    near-zero investment and health); the matched composition battery and the
    transmission sign-flip re-run are pending on the current calibration ---
    quote neither until they land.}
\end{frame}

\figframe{Who gains and who loses is not a budget-share story}%
  {impact/impact_incidence}%
  {Consumption-equivalent incidence by lifetime-income group on the coupled
   run: middle groups gain, the 90th--99th percentile group loses most, and
   baseline energy budget shares are nearly uniform --- the dispersion runs
   through earnings, assets, and transfers. PEP vs Base, Philippines.}
